\section{Conventions, notation, and variations of standard facts}

\subsection{Global conventions}

Throughout this paper, all schemes, varieties and morphisms will be defined over
the complex number field. We follow the notation and conventions of
Hartshorne's book \cite{Ha77}. In particular, varieties are always assumed to
be irreducible. For all notation around Mori theory, such as klt spaces and klt
pairs, we refer the reader to \cite{KM98}.

\subsection{Varieties and complex spaces}

In order to keep notation simple, we will sometimes, when there is no danger of
confusion, not distinguish between algebraic varieties and their underlying
complex spaces. Along these lines, if $X$ is a quasi-projective complex
variety, we write $\pi_1(X)$ for the fundamental group of the associated complex
space.

\subsection{Reflexive sheaves}

As in most other papers on the subject, we will frequently consider reflexive
sheaves and take reflexive hulls. Given a normal, quasi-projective variety (or
normal, irreducible complex space) $X$, we write
$\Omega^{[p]}_X := \bigl(\Omega^p_X \bigr)^{**}$ and refer to this sheaf as \emph{the
sheaf of reflexive differentials}. More generally, given any coherent sheaf
$\sE$ on $X$, write $\sE^{[\otimes m]} := \bigl(\sE^{\otimes m} \bigr)^{**}$ and
$\det \sE := \bigl(\wedge^{\rank \sE} \sE \bigr)^{**}$. Given any morphism $f : Y \to X$ of
normal, quasi-projective varieties (or normal, irreducible, complex spaces), we
write $f^{[*]} \sE := (f^* \sE)^{**}$.

